% Generated draft: requires mathematical and source review.
\documentclass[11pt]{article}
\usepackage[margin=25mm]{geometry}
\usepackage{fontspec}
\setmainfont{Latin Modern Roman}
\usepackage{amsmath,amssymb,mathtools,unicode-math}
\setmathfont{Latin Modern Math}
\usepackage{enumitem,needspace,xurl,hyperref}
\setlength{\parindent}{0pt}
\setlength{\parskip}{5pt}
\setlength{\emergencystretch}{3em}
\newcommand{\N}{\mathbb{N}}
\newcommand{\Z}{\mathbb{Z}}
\newcommand{\Q}{\mathbb{Q}}
\newcommand{\R}{\mathbb{R}}
\newcommand{\C}{\mathbb{C}}
\newcommand{\F}{\mathbb{F}}
\newcommand{\A}{\mathbb{A}}
\newcommand{\PP}{\mathbb{P}}
\newcommand{\E}{\mathbb{E}}
\newcommand{\eps}{\varepsilon}
\newcommand{\dd}{\,\mathrm{d}}
\begin{document}
\section*{UnsolvedMath ID: 30003858}
\textbf{Source problem code:} catalogue entry
\subsection{Definitions and mathematical statement}
The source leaves the dimension range, target geometry, regularity, and initial-trace convention unspecified. The following definitions describe the customary Euclidean interpretation, conditional on equipping the target \(N\) with a Riemannian metric.

Given the angular datum \(u_0:S^{n-1}\to N\), its \(0\)-homogeneous extension is
\[
\bar u_0:\R^n\setminus\{0\}\to N,\qquad
\bar u_0(x)=u_0\!\left(\frac{x}{|x|}\right).
\]
Thus \(\bar u_0(\lambda x)=\bar u_0(x)\) for every \(x\ne0\) and \(\lambda>0\). No value at the origin is specified.

For a smooth map \(u:\R^n\times(0,\infty)\to N\), the harmonic-map heat-flow equation is
\[
\partial_tu=\tau(u),\qquad
\tau(u)=\operatorname{tr}_{\R^n}\nabla\dd u,
\]
where the tension field uses the Euclidean metric on the domain and the Riemannian metric on \(N\). Such a solution is an expanding self-similar solution, or expander, if
\[
u(\lambda x,\lambda^2t)=u(x,t)
\quad\text{for every }x\in\R^n,\ t>0,\ \lambda>0.
\]
Equivalently, with \(U(y)=u(y,1)\),
\[
u(x,t)=U\!\left(\frac{x}{\sqrt t}\right),
\qquad
\tau(U)(y)+\frac12\dd U_y(y)=0.
\]
These classical definitions do not specify the weak solution classes mentioned in the supplied literature assessment.

The requested problem is to characterize, among harmonic-map-flow solutions emerging from \(\bar u_0\), those that are expanders, and to determine the hypotheses, admissibility conditions, and meaning of uniqueness under which an appropriate uniqueness theorem for expanding solutions can be formulated and proved. Here ``emerging from'' requires a specified initial-trace notion, which the source does not provide. The question therefore remains a formulation problem as well as a request for a detection criterion and a uniqueness theorem; it does not supply a precise assertion of unconditional uniqueness.

The supplied dated curation reports existence for homotopically trivial Lipschitz data, generic uniqueness through relative entropy, and nonuniqueness in a hemisphere setting in work dated 2026. The applicable dimensions, target assumptions, solution classes, and meaning of genericity are not supplied. These qualifications must be retained when formulating a theorem; they are not independently verified results in this draft.
\subsection{Short English statement}
Given sphere-valued angular data with values in a target manifold, extend the data constantly along rays and consider harmonic-map heat flows with that initial trace. Determine how to recognize flows that expand by parabolic scaling and under what precise hypotheses their uniqueness should be expected and proved. The source leaves the target assumptions, solution class, trace convention, and uniqueness conditions unresolved. Its partially solved label and literature claims are supplied dated curation.
\subsection{Sources}
\begin{itemize}
\item Local UnsolvedMath registry, supplied complete source record, permanent ID \(30003858\), title ``Detection and Uniqueness of Expanding Harmonic-Map Flows,'' exact source problem code \texttt{catalogue entry}, collection ``Oberwolfach Reports Open Problems.'' The record supplies proposed year 2018 but names no proposer.
\item Supplied report citation: \emph{Geometrie (2019)}, Oberwolfach Reports, pp.~1635--1684, DOI \texttt{10.4171/owr/2018/27}. \url{https://doi.org/10.4171/owr/2018/27}. The background describes Oberwolfach Reports as documenting mathematical workshops and their open-problem discussions. The citation is supplied provenance and has not been independently checked.
\item Supplied literature reference, retaining its exact label: ``Deruelle--Lamm, Expanding Ricci solitons.'' \url{https://arxiv.org/abs/1801.08012}. The label concerns Ricci solitons, whereas the problem concerns harmonic-map flow; its bibliographic identity and relevance require verification.
\item Supplied literature reference: ``Deruelle, Relative entropy and uniqueness.'' \url{https://doi.org/10.1080/03605302.2019.1646280}.
\item Supplied literature reference: ``Li, Nonuniqueness of weak expanders.'' \url{https://arxiv.org/abs/2602.08932}.
\item Literature assessment dated exactly \texttt{2026-08-21}: the record labels the problem partially solved and reports existence for homotopically trivial Lipschitz data, generic uniqueness through relative entropy, unsettled general detection and uniqueness questions, and 2026 nonuniqueness in a hemisphere setting. This is supplied dated curation derived from the OWR record, not an independently verified literature assessment; the record itself requests independent verification of time-sensitive status and cited claims.
\item The supplied statement status is \texttt{corrected\_verified}. Its accompanying explanation describes repairs to notation, typography, and sentence structure while preserving mathematical content. This editorial status does not establish mathematical correctness or completeness of the formulation.
\end{itemize}
\end{document}
